\documentclass[acmsmall,screen,review,anonymous]{acmart}

\usepackage{amsmath,mathtools}
\usepackage{booktabs}
\usepackage{multirow}
\usepackage{array}
\usepackage{tabularx}
\usepackage{graphicx}
\usepackage{xcolor}
\usepackage{microtype}
\usepackage{balance}

\settopmatter{printfolios=true,printacmref=false}
\setcopyright{none}
\renewcommand\footnotetextcopyrightpermission[1]{}
\acmDOI{}
\acmJournal{TAAS}
\acmYear{2026}
\acmVolume{1}
\acmNumber{1}
\acmArticle{1}
\acmMonth{1}
\citestyle{acmauthoryear}

\newcommand{\rrmse}{\mathrm{RRMSE}}
\newcommand{\R}{\mathbb{R}}
\newcolumntype{Y}{>{\raggedright\arraybackslash}X}
\newcommand{\figureplaceholder}[2]{%
  \fbox{\begin{minipage}[c][#1][c]{0.96\linewidth}
  \small\textbf{Figure design placeholder.} #2
  \end{minipage}}}

\title[Subject-Calibrated EOG-Guided Diffusion]{Subject-Calibrated EOG-Guided Diffusion for EEG Denoising}

\author{Anonymous Authors}
\affiliation{%
  \institution{Anonymous Institution}
  \city{Anonymous}
  \country{Anonymous}}
\email{anonymous@example.org}

\begin{abstract}
Ocular artifacts vary across participants because the coupling between eye activity and scalp electrodes depends on anatomy, electrode placement, and recording conditions. Population denoisers learn common electroencephalographic (EEG) structure but do not explicitly model this participant-dependent contamination pathway. We present a subject-calibrated, electrooculography (EOG)-guided diffusion method that separates neural-signal modeling from ocular-artifact calibration. A population-trained conditional diffusion model is shared across participants, while a temporally separate calibration segment estimates an EOG-to-EEG propagation matrix. Empirical Bayes shrinkage stabilizes this estimate and replaces unreliable individual deviations with population calibration. In same-session paired semi-simulation with 15 development participants, matched EOG guidance reduced participant-level relative root-mean-square error by 0.1428 relative to a zero-anchor control that retained the same calibration state, with improvement for all participants. Evaluation on eight held-out participants produced a similar improvement of 0.1537 (95\% bootstrap interval 0.0725--0.2504; 7/8 participants). On natural development recordings without a clean target, matched calibration increased EOG attenuation by 1.55 dB and low-EOG signal retention by 0.071 relative to population calibration. Propagation-aware sample dispersion yielded empirical 80\% and 90\% interval coverages of 0.802 and 0.853 on development paired data, with continuous ranked probability score 0.1503. These results show that a short, query-disjoint EEG--EOG calibration can adapt a shared diffusion denoiser to participant-specific ocular propagation.
\end{abstract}

\ccsdesc[500]{Computing methodologies~Machine learning}
\ccsdesc[500]{Applied computing~Health informatics}
\ccsdesc[300]{Human-centered computing~Ubiquitous and mobile computing systems and tools}

\keywords{EEG denoising, ocular artifacts, electrooculography, conditional diffusion, subject calibration, predictive intervals}

\begin{document}
\maketitle

\begin{figure*}[t]
  \centering
  \figureplaceholder{34mm}{\textbf{Graphical abstract.} A compact left-to-right overview with four visual blocks: (1) a short, temporally separate EEG--EOG calibration segment; (2) estimation and empirical Bayes shrinkage of the participant's EOG-to-EEG propagation matrix; (3) evaluation EEG--EOG and the matched propagation guide entering the shared conditional diffusion model, with population calibration and a zero-anchor control shown as comparisons; and (4) restored EEG and sample-based predictive intervals. Use blue for the shared population model, orange for calibration-derived quantities, green for matched guidance, and gray for comparison modes. This panel is conceptual and contains no synthetic EEG traces or performance values.}
  \caption{Graphical abstract design. The method adapts the ocular-artifact observation model from a separate calibration segment while keeping the diffusion denoiser shared across participants.}
  \Description{Placeholder specifying a four-stage graphical abstract for subject-calibrated EOG-guided EEG denoising.}
  \label{fig:graphical-abstract}
\end{figure*}

\section{Introduction}
\label{sec:introduction}

Electroencephalography (EEG) records neural activity at the scalp with millisecond resolution, but the same electrodes also capture potentials generated by blinks and eye movements. These ocular artifacts are often larger than the EEG activity of interest and spread across many channels. Rejecting every affected segment can remove a substantial fraction of a recording, while aggressive correction can attenuate event-related potentials, oscillatory activity, or spatial structure. Ocular-artifact removal is therefore a restoration problem: the method must reduce the contribution of the eyes while preserving the neural signal that overlaps it in time and frequency.

Existing approaches address this problem through EOG regression, adaptive filtering, blind source separation, artifact-subspace methods, and supervised neural networks \cite{gratton1983ocular,semlitsch1986ocular,jung2000bss,chang2020asr,zhang2021eegdenoisenet,yu2022deepseparator}. Classical regression is attractive when electrooculography (EOG) is recorded because its correction coefficients have a direct interpretation, but fixed coefficients do not capture changes in electrode placement and ocular propagation. Learned denoisers provide more flexible mappings, yet a single population model must average over these participant- and session-dependent differences. Diffusion models offer a conditional generative formulation for EEG restoration \cite{huang2025eegdfus}, but conditioning on the corrupted waveform alone does not identify how a particular participant's eye activity reaches the scalp electrodes.

The key observation behind this work is that the participant-dependent quantity needed for ocular correction is not an identity label. It is the EOG-to-EEG propagation measured during the recording session. Calibration-based ocular correction has long estimated this propagation from simultaneous EEG and EOG \cite{gratton1983ocular,kobler2020sgeyesub}; we use the same physical cue to adapt a population conditional diffusion model. A short calibration segment is kept separate from evaluation data and used to estimate a propagation matrix. At evaluation time, the recorded EOG and this matrix produce an ocular-artifact estimate that guides the denoiser. Empirical Bayes shrinkage stabilizes the estimate and returns low-reliability calibration to the population estimate rather than forcing an individual correction.

This design yields a single subject-calibrated EOG-guided diffusion method. The neural-signal model is shared across participants, whereas the ocular observation model is calibrated within each participant--session--task cell. The separation permits direct controls: population calibration tests whether an individual estimate is useful, calibration from another participant tests specificity, and a zero-anchor control measures the contribution of the EOG-derived waveform guide while retaining the same network and compact calibration state. Repeated diffusion samples are also used to form empirical predictive intervals, with part of their width assigned to uncertainty in the estimated propagation matrix.

We evaluate the method on a development cohort and a held-out participant cohort, using paired semi-simulation for waveform error and natural recordings for EOG attenuation and signal-preservation measures. Matched calibration improves over both population calibration and the zero-anchor control in development, and the zero-anchor comparison is repeated without method selection in held-out participants. A controlled injected sensitivity analysis examines whether a support-fitted representation can carry the same participant-matched calibration information through a different interface. Calibration-aware predictive intervals improve empirical coverage without increasing the continuous ranked probability score relative to the unadjusted sample distribution. The experiments support a practical separation of roles: a population-trained diffusion model represents common EEG structure, while a short EEG--EOG calibration adapts the artifact pathway.

\section{Related Work}
\label{sec:related-work}

Ocular-artifact correction has traditionally relied on an auxiliary EOG recording or on spatial separation within the EEG. Regression methods estimate how ocular potentials propagate to each scalp channel and subtract the predicted contribution \cite{gratton1983ocular,semlitsch1986ocular}. Frequency-domain regression and adaptive filtering extend this view to frequency-dependent or time-varying coupling, whereas independent component analysis and artifact subspace reconstruction separate artifacts from neural activity without an explicit EOG-to-EEG transfer model \cite{jung2000bss,chang2020asr}. Calibration-based subspace subtraction has shown that eye-related components can be estimated once and applied online through matrix operations \cite{kobler2020sgeyesub}. These methods establish the value of a calibration segment, but they do not combine participant-specific ocular propagation with a learned distribution of clean EEG.

Deep EEG denoisers learn nonlinear mappings from contaminated to clean or cleaner signals. EEGdenoiseNet introduced a widely used semi-simulated benchmark, and subsequent convolutional, decomposition-based, and U-Net models improved artifact removal under related paired constructions \cite{zhang2021eegdenoisenet,sun2020rescnn,yu2022deepseparator,chuang2022icunet,lai2024cleegn}. Their main strength is flexible waveform reconstruction; their usual limitation for the present setting is that participant variation is absorbed into shared model parameters. A participant label or a generic recording embedding does not directly specify the ocular component to remove. Our method instead retains a shared denoiser and conditions it on a propagation estimate obtained from simultaneous EEG and EOG during a separate calibration interval.

Diffusion models learn a reverse process from a prescribed noising process and can be conditioned on degraded observations for image restoration \cite{ho2020ddpm,song2021scoresde,saharia2022palette}. Methods such as DDRM and diffusion posterior sampling incorporate a known degradation process during reconstruction \cite{kawar2022ddrm,chung2023dps}; EEGDfus applies conditional diffusion directly to EEG denoising \cite{huang2025eegdfus}. Our formulation is closest in spirit to conditional restoration, but the condition is factored into a corrupted EEG observation and a separately calibrated EOG-derived artifact estimate. The method also returns sample-based predictive intervals. We evaluate these intervals through empirical marginal coverage and the continuous ranked probability score rather than treating stochastic variation as calibrated uncertainty by construction \cite{gneiting2007strictly}.

\section{Methodology}
\label{sec:methodology}

\subsection{Calibration of EOG-to-EEG propagation}

For participant--session--task cell $i$, let $\mathbf{y}\in\R^{C\times T}$ denote observed EEG, $\mathbf{x}$ the latent neural signal, and $\mathbf{e}\in\R^{R\times T}$ the simultaneously recorded EOG. Over a calibration and evaluation interval, we use the local additive model
\begin{equation}
  \mathbf{y}=\mathbf{x}+\mathbf{A}_i\mathbf{e}+\boldsymbol{\varepsilon},
  \label{eq:observation}
\end{equation}
where $\mathbf{A}_i\in\R^{C\times R}$ is the EOG-to-EEG propagation matrix and $\boldsymbol{\varepsilon}$ contains residual activity not represented by the ocular model. Equation~\ref{eq:observation} is an operational description of ocular propagation over the recording interval, rather than a complete physiological forward model.

The calibration segment contains paired EEG and EOG and is temporally disjoint from all evaluation segments. Each EOG channel is centered by its calibration median and scaled by $1.4826$ times its median absolute deviation, with a $10^{-6}$ lower bound; EEG uses channel scales estimated from the training fold. After temporal centering, ridge regression produces
\begin{equation}
  \widehat{\mathbf{A}}_i=\mathbf{Y}_c\mathbf{E}_c^\top
  \left(\mathbf{E}_c\mathbf{E}_c^\top+r\mathbf{I}\right)^{-1},
  \qquad
  r=0.05\,\max\!\left\{
  \frac{\operatorname{tr}(\mathbf{E}_c\mathbf{E}_c^\top)}{R},\epsilon\right\}.
  \label{eq:ridge}
\end{equation}
where $\epsilon$ is machine precision for the numerical implementation.
The recipient-excluded population matrix $\mathbf{A}_{\mathrm{pop}}$ is the mean of training-participant ridge estimates from the first 30 s of the same session--task cell. To reduce variance in the 120-s individual estimate, it is shrunk toward this population matrix,
\begin{equation}
  \widetilde{\mathbf{A}}_i=
  \mathbf{A}_{\mathrm{pop}}+\lambda_i
  (\widehat{\mathbf{A}}_i-\mathbf{A}_{\mathrm{pop}}),
  \qquad 0\leq\lambda_i\leq1.
  \label{eq:shrinkage}
\end{equation}
The 120-s calibration segment is divided into four consecutive blocks. Let the scalar $\tau^2$ be the mean squared deviation, across training participants and matrix entries, of their full 120-s estimates from $\mathbf{A}_{\mathrm{pop}}$; let the scalar $v_i$ be the corresponding mean squared deviation of the four 30-s block estimates from participant $i$'s full estimate. We set $\lambda_i=\tau^2/(\tau^2+v_i/4)$, clipped to $[0,1]$. If less than 60 s is available or $v_i$ exceeds the 95th percentile estimated in the training fold, $\lambda_i$ is set to zero, reducing $\widetilde{\mathbf{A}}_i$ exactly to $\mathbf{A}_{\mathrm{pop}}$. This computation uses calibration EEG and EOG but no clean evaluation target or participant identifier.

\subsection{Conditional diffusion restoration}

At evaluation time, recorded EOG is converted into an ocular waveform guide. For mode $m$, the guide is
\begin{equation}
  \mathbf{a}^{(m)}_i=
  \begin{cases}
    \widetilde{\mathbf{A}}_i\mathbf{e}, & m=\text{matched},\\
    \mathbf{A}_{\mathrm{pop}}\mathbf{e}, & m=\text{population},\\
    \mathbf{0}, & m=\text{zero anchor}.
  \end{cases}
  \label{eq:guide-modes}
\end{equation}
The zero-anchor mode retains the matched compact calibration state, so it isolates the explicit waveform guide rather than removing every calibration-derived input. It is an ablation, distinct from the reliability rule above: unreliable individual calibration uses the population matrix and population state.

The EEG model is shared across participants. A population-trained one-dimensional U-Net receives the diffusion state $\mathbf{x}_t$, contaminated EEG $\mathbf{y}$, guide $\mathbf{a}^{(m)}_i$, diffusion time $t$, and a mode-specific calibration state $\mathbf{c}^{(m)}_i\in\R^{C\times 53}$. For each EEG channel, this state contains the two propagation coefficients, their log norm, four calibration-quality variables (vertical and horizontal EOG magnitude, ridge fit, and condition number), and a 46-dimensional one-hot channel indicator. A row encoder and one-dimensional spatial convolution reduce this state to a 128-dimensional context and channelwise modulation weights. Matched and zero-anchor modes use the reliability-gated matched state, which equals the population state when the calibration is rejected; population calibration always uses the population state. The clean estimate is parameterized around EOG regression,
\begin{equation}
 \widehat{\mathbf{x}}_0=
 (\mathbf{y}-\mathbf{a}^{(m)}_i)
 +\Delta_{\mathrm{pop},\theta}(\mathbf{x}_t,\mathbf{y},\mathbf{a}^{(m)}_i,t)
 +\Delta_{\mathrm{cal},\theta}(\mathbf{x}_t,\mathbf{y},\mathbf{a}^{(m)}_i,t,
 \mathbf{c}^{(m)}_i).
 \label{eq:clean-estimate}
\end{equation}
The first residual term in Equation~\ref{eq:clean-estimate} models EEG structure shared across participants; the second permits the correction to depend on the calibrated propagation. Both are learned jointly by minimizing Smooth L1 loss between the predicted and reference clean EEG. We use deterministic DDIM updates for sampling \cite{song2020ddim}. The role of the diffusion model is to represent a population distribution of plausible clean waveforms conditioned on the observation and its calibrated ocular component.

The same network parameters are used for every mode in Equation~\ref{eq:guide-modes}; the method does not retrain a participant-specific denoiser. Figure~\ref{fig:method} summarizes the calibration and evaluation flows.

\begin{figure*}[t]
  \centering
  \figureplaceholder{48mm}{\textbf{Method overview.} Panel A should show non-overlapping calibration and evaluation timelines. Panel B should show simultaneous calibration EEG and EOG entering median/MAD scaling, ridge regression, and empirical Bayes shrinkage to estimate $\widetilde{\mathbf{A}}_i$, with low-reliability estimates reverting to $\mathbf{A}_{\mathrm{pop}}$. Panel C should show evaluation EEG and EOG forming the matched guide $\widetilde{\mathbf{A}}_i\mathbf{e}$. Panel D should show the shared one-dimensional U-Net inside the DDIM reverse process, with population calibration and the zero-anchor control distinguished from the proposed matched mode. Solid arrows denote evaluation-time inputs; dashed arrows denote calibration-only computations. Do not include synthetic performance values or invented waveforms.}
  \caption{Design of the subject-calibrated EOG-guided diffusion method. Participant adaptation occurs through the EOG-to-EEG propagation matrix; the conditional diffusion model is shared across participants.}
  \Description{Placeholder for a four-panel diagram of calibration, propagation estimation, reliability-based routing, and conditional diffusion restoration.}
  \label{fig:method}
\end{figure*}

\subsection{Sample-based predictive intervals}

For uncertainty analysis, $K=32$ reverse trajectories jointly sample DDIM initial noise and an entrywise propagation matrix,
\begin{equation}
  A^{(k)}_{i,cr}\sim
  \mathcal{N}\!\left(\widetilde A_{i,cr},\sigma^2_{A_i,cr}\right).
  \label{eq:operator-sample}
\end{equation}
The entrywise variance uses the inverse-variance combination
\begin{equation}
  \sigma^2_{A_i,cr}=\left(\frac{1}{\tau^2_{cr}}+\frac{4}{v_{i,cr}}\right)^{-1},
  \qquad
  w^2_{A_i,c}(t)=\sum_{r=1}^{R}\sigma^2_{A_i,cr}e_r(t)^2,
  \label{eq:operator-variance}
\end{equation}
where $\tau^2_{cr}$ is the between-participant variance and $v_{i,cr}$ is the four-block within-calibration variance for matrix entry $(c,r)$. The trajectory mean is used as the point estimate. The empirical trajectory width $w_{\mathrm{chain}}$ is combined with the propagated EOG-to-EEG variance from Equation~\ref{eq:operator-variance},
\begin{equation}
  w^2=w_{\mathrm{chain}}^2+w_A^2,
  \qquad
  I_{1-\alpha}=\overline{\mathbf{x}}
  \mathbin{\pm} \tau_{-f}z_{1-\alpha/2}w,
  \label{eq:interval}
\end{equation}
where $\tau_{-f}$ is a scalar temperature selected without evaluation fold $f$, and $z_p$ is the standard-normal quantile at probability $p$. The smallest temperature on the grid 0.50--6.00 in increments of 0.05 that attained 80\% coverage in the remaining folds was used. We refer to $I_{1-\alpha}$ in Equation~\ref{eq:interval} as a sample-based predictive interval. Its quality is measured empirically through marginal coverage, interval width, continuous ranked probability score (CRPS), and risk--coverage curves.

\section{Experiments}
\label{sec:experiments}

\subsection{Data and evaluation protocol}

The main experiments used multimodal recordings with synchronized EEG and EOG from a mobile brain--computer interface dataset \cite{guttmannflury2025ocular,guttmannflury2024eyebci}. Signals were represented by 46 EEG channels and two bipolar EOG derivations (vertical and horizontal), resampled to 100 Hz. Each model input contained 512 samples (5.12 s). Within each available session--task cell, a continuous 120-s prefix supplied the participant calibration, and all evaluation windows occurred later in the record. Fifteen participants were assigned to five participant-held-out folds, each with nine training, three validation, and three test participants. Eight additional participants were reserved for a held-out evaluation after the model and analysis had been selected. Clean targets, task outcomes, and evaluation labels were never model inputs. The source study was approved by the Institutional Review Board of Shanghai Jiao Tong University (IRB HRP E2021216I), and participants provided written consent for anonymized data use; the present work is a secondary analysis of the public, de-identified recordings \cite{guttmannflury2025ocular}.

Waveform accuracy was measured with paired semi-simulation. A low-ocular-activity EEG carrier served as the reference, and a separately sampled ocular waveform was passed through the recipient's later-record propagation matrix to produce a contaminated observation. That generator-only matrix was estimated from 150--270 s and never entered the model. EEG carriers and ocular waveforms were sampled from 5.12-s windows beginning at or after 300 s; draws could overlap across episodes. The EEG carrier, ocular waveform, and recipient calibration came from different participants, isolating the propagation component while retaining real EEG and EOG signals. For nonzero-artifact episodes, the injected amplitude used one of three base gains (0.35, 0.70, or 1.15) multiplied by a uniform factor from 0.85 to 1.15; 15\% of sampled episodes contained exactly zero injected artifact. Eight paired episodes were evaluated per participant and distributed across available session--task cells. Natural recordings were analyzed separately because they have no clean target; four windows were placed linearly across the interval beginning at 300 s in each available cell. Their endpoints were EOG attenuation, EOG--EEG coherence reduction, and low-EOG signal retention; paired and natural outcomes were never combined into one score.

An auxiliary calibration-representation experiment used EEGEyeNet \cite{kastrati2021eegeyenet}. A compact deterministic model received a 46-channel, 5.12-s injected episode at 100 Hz and either a participant-matched, zero-embedding, or mismatched calibration representation; each participant's episodes were divided chronologically into calibration and query halves. For each training-pool size from 30 to 259 participants, the model was trained for 20,000 AdamW updates (batch 32, learning rate $2\times10^{-4}$, EMA 0.999) with artifact-reconstruction mean squared error. With the network then fixed, a 32-dimensional representation for each evaluation participant was optimized for 200 steps against paired clean calibration targets and evaluated on the disjoint query episodes. This is an evaluator-only upper-bound analysis, not an input to the main method. The complete 15-participant result is reported first. A subsequent plausibility sensitivity analysis restricted the injected-artifact-to-clean-signal ratio to $[0.1,20]$ and retained 14 participants. Table~\ref{tab:panels} keeps this controlled study separate from the denoising cohorts.

\begin{table*}[t]
\caption{Experimental panels. Sample sizes refer to independent participants, not windows or diffusion samples.}
\label{tab:panels}
\centering
\small
\begin{tabular}{@{}p{.22\textwidth}p{.13\textwidth}p{.24\textwidth}p{.30\textwidth}@{}}
\toprule
Panel & Participants & Inputs & Primary purpose \\
\midrule
Development denoising & 15 & Evaluation EEG and EOG; separate calibration segment & Participant-held-out method comparison, natural-signal measures, and predictive intervals \\
Held-out denoising & 8 & Frozen five-fold model ensemble; same paired construction & Held-out evaluation of matched calibration against the zero-anchor control \\
Calibration representation & 259 train; 15 evaluation; 14 in plausible range & Injected episode and calibration-derived condition & Matched, zero-embedding, and mismatched conditioning comparison \\
\bottomrule
\end{tabular}
\end{table*}

\subsection{Comparison methods and outcome measures}

All calibration variants used the same population diffusion network and the same evaluation instances. \emph{Matched calibration} used the participant's own propagation estimate after shrinkage. \emph{Population calibration} used the recipient-excluded pooled estimate. \emph{Mismatched calibration} used a fixed eligible participant from the same session--task stratum. The \emph{zero-anchor control} set the ocular waveform guide to zero while retaining the matched compact calibration state, whereas \emph{shuffled EOG} disrupted the temporal alignment between EOG and EEG. Raw EEG, linear EOG regression, and a one-step deterministic network were included when they shared the corresponding evaluation protocol. Table~\ref{tab:variants} states the information available to each comparison.

The primary paired endpoint was participant-level RRMSE, where lower values are better. We report a utility $U(A>B)=\rrmse(B)-\rrmse(A)$, so positive values favor method $A$. The development comparison used matched calibration against both population calibration and the zero-anchor control; the held-out comparison used matched calibration against the same zero-anchor control. Natural-signal endpoints were reported individually rather than collapsed into a composite. Predictive intervals were evaluated at nominal 50\%, 80\%, and 90\% levels with empirical marginal coverage, Gaussian CRPS, and risk--coverage area.

\begin{table*}[t]
\caption{Method variants and information available during evaluation. All learned variants use the same population model unless otherwise indicated.}
\label{tab:variants}
\centering
\small
\begin{tabularx}{\textwidth}{@{}p{.17\textwidth}Yp{.18\textwidth}p{.13\textwidth}p{.19\textwidth}@{}}
\toprule
Variant & Propagation source & Use in the model & Evaluation EOG & Role \\
\midrule
Matched calibration & Same participant-session & Waveform guide and compact state & Recorded & Proposed method \\
Population calibration & Recipient-excluded pool & Waveform guide and compact state & Recorded & Pooled-calibration control \\
Mismatched calibration & Another participant & Waveform guide and compact state & Recorded & Specificity control \\
Zero-anchor control & Same participant-session & Compact state only; waveform guide set to zero & Not used in guide & Explicit-guide ablation \\
Shuffled EOG & Same participant-session & Misaligned guide and matched compact state & Temporally shuffled & Temporal-alignment control \\
Linear EOG regression & Same participant-session & Direct linear subtraction & Recorded & Classical reference \\
Deterministic network & Same inputs in its protocol & One-step reconstruction & Recorded & Learned point-estimate reference \\
\bottomrule
\end{tabularx}
\end{table*}

\subsection{Training and statistical analysis}

Population matrices, normalization statistics, shrinkage quantities, model parameters, and reliability thresholds were fitted without the evaluated participant. Throughout the point-estimate, held-out, and interval analyses reported here, a rejected calibration reverted to the population propagation matrix and population state. The four-scale one-dimensional U-Net used channel widths 32, 64, 96, and 128, three downsampling stages, four-head self-attention at the bottleneck, and a 128-dimensional calibration context; the instantiated architecture had 965,085 trainable parameters. It predicted clean EEG directly over 1,000 linear diffusion levels ($\beta_1=10^{-4}$, $\beta_{1000}=0.02$). Each fold and seed was trained for 80,000 AdamW updates (batch 8, learning rate $10^{-4}$, weight decay $10^{-4}$), using Smooth L1 loss, gradient clipping at 1.0, and exponential moving average 0.999. During training, the waveform anchor was dropped with probability 0.30 and the compact state was replaced by its population counterpart with probability 0.20. Participant-aggregated population RRMSE on the validation set, evaluated every 2,000 updates, selected the moving-average parameters.

Development point estimates used five folds, three fixed training seeds (20261201--20261203), 50 deterministic DDIM steps, and one trajectory per fold--seed cell; initial-noise draws were identical across calibration variants. The held-out estimate used the seed-20261201 model from each fold, a population registry built from all 15 development participants, and an output-space mean of the five fold predictions with a common initial-noise draw for each episode. It is therefore a frozen model average rather than five stochastic trajectories. Only the uncertainty analysis used 32 stochastic trajectories per development fold--seed cell, with interval temperature estimated while leaving out the evaluated fold.

For point-estimate, natural-signal, and calibration-representation contrasts, windows, repeated protocols, seeds, and calibration conditions were first averaged within participant. We then report the participant-level mean, median, number of positive effects, and descriptive 95\% percentile bootstrap intervals from 5,000 participant resamples. For predictive intervals, coverage was computed over channel--time elements within each fold--seed evaluation cell and then averaged over the 15 cells; Gaussian CRPS and risk--coverage area were aggregated over evaluation episodes. Neither signal elements nor episodes were treated as independent participants for inferential claims.

\section{Results}
\label{sec:results}

\subsection{Matched calibration improves paired EEG restoration}

On the 15-participant development cohort, matched calibration achieved participant-level RRMSE 0.43097, compared with 0.57378 for the zero-anchor control and 0.65114 for population calibration (Table~\ref{tab:main-results}). The improvement over zero anchor was 0.14281 (95\% bootstrap interval 0.10613--0.18342) and was positive for all 15 participants, identifying the contribution of the explicit matched waveform guide beyond the shared network and compact calibration state. The improvement over population calibration was 0.22017 (0.09293--0.40245) and was positive for 14 of 15 participants. Thus, neither suppressing the waveform anchor nor replacing the individual propagation estimate by a pooled estimate reproduced the matched result.

The same method and zero-anchor comparison were then evaluated on eight held-out participants. Matched calibration achieved RRMSE 0.53509, whereas the zero-anchor control achieved 0.68882. The participant-level improvement was 0.15373 (0.07245--0.25039), with a median of 0.14740 and positive effects for 7 of 8 participants. Its magnitude was close to the development effect despite the higher absolute error in the held-out cohort. Figure~\ref{fig:main-results} is reserved for the paired waveforms and participant-level contrasts from these two cohorts.

\begin{table*}[t]
\caption{Participant-level paired restoration results. Lower RRMSE is better. Development and held-out cohorts are reported separately and are not pooled.}
\label{tab:main-results}
\centering
\small
\begin{tabularx}{\textwidth}{@{}lYcYYc@{}}
\toprule
Cohort & Variant & RRMSE & Primary comparison & Mean improvement (95\% interval) & Positive participants \\
\midrule
\multirow{4}{*}{Development ($n=15$)}
 & Raw EEG & 0.60764 & -- & -- & -- \\
 & Population calibration & 0.65114 & Matched vs population & 0.22017 (0.09293--0.40245) & 14/15 \\
 & Zero anchor & 0.57378 & Matched vs zero anchor & 0.14281 (0.10613--0.18342) & 15/15 \\
 & Matched calibration & \textbf{0.43097} & -- & -- & -- \\
\midrule
\multirow{4}{*}{Held-out ($n=8$)}
 & Raw EEG & 0.77531 & -- & -- & -- \\
 & Population calibration & 0.76607 & -- & -- & -- \\
 & Zero anchor & 0.68882 & Matched vs zero anchor & 0.15373 (0.07245--0.25039) & 7/8 \\
 & Matched calibration & \textbf{0.53509} & -- & -- & -- \\
\bottomrule
\end{tabularx}
\end{table*}

\begin{figure*}[t]
  \centering
  \figureplaceholder{58mm}{\textbf{Paired restoration and participant effects.} Panel A should display recorded EOG above aligned raw EEG, paired reference EEG, linear EOG regression, zero-anchor diffusion, and matched-calibration diffusion for three fixed scalp channels. Select the lowest-identifier development participant with complete outputs, the earliest eligible ocular event, and the same channel names and vertical scale for every variant; report these identifiers in the final caption. Panel B should show participant-paired RRMSE changes from zero anchor to matched calibration for all 15 development participants. Panel C should show the corresponding paired changes for all eight held-out participants. Panel D should show participant-level effect estimates with bootstrap intervals for the two cohorts on separate axes. No waveform or participant point may be synthesized for illustration.}
  \caption{Planned visualization of paired denoising and participant-level effects. Waveforms will be drawn from stored evaluator outputs using the deterministic selection rule stated in the panel description; lower RRMSE indicates better reconstruction.}
  \Description{Placeholder for real EOG and EEG waveform comparisons plus paired participant plots for development and held-out cohorts.}
  \label{fig:main-results}
\end{figure*}

The matched diffusion estimate was comparable to the classical linear EOG reference in the development protocol: RRMSE was 0.43097 for matched diffusion, 0.43527 for linear regression, and 0.49682 for the one-step deterministic EOG network. These comparisons locate the gain in participant-calibrated EOG guidance rather than in a general point-error advantage of diffusion. In a separate retraining experiment that supplied population calibration to both learned estimators, RRMSE was 0.50430 for the deterministic network and 0.52598 for diffusion. The diffusion sampler is retained here for conditional waveform modeling and predictive sampling; linear and deterministic estimators remain strong point-estimate references.

\subsection{Natural recordings and calibration controls}

On natural development recordings, matched calibration achieved 2.463 dB EOG attenuation and low-EOG signal retention of 0.843. Population calibration achieved 0.909 dB and 0.773, respectively, while the zero-anchor control preserved 0.921 of the low-EOG signal but attenuated only 0.277 dB (Table~\ref{tab:natural-results}). Relative to population calibration, matched calibration increased attenuation by 1.554 dB (0.969--2.211; 14/15 participants), retention by 0.0707 (0.0382--0.1067; 13/15), and EOG--EEG coherence reduction by 0.0849 (0.0535--0.1211; 14/15). The joint improvement indicates that the paired-error gain was not obtained by simply suppressing a larger fraction of every evaluation segment.

\begin{table}[t]
\caption{Natural-recording operating points on the development cohort. Natural recordings have no clean target. Higher attenuation and retention are preferred.}
\label{tab:natural-results}
\centering
\small
\begin{tabular}{@{}lcc@{}}
\toprule
Variant & EOG attenuation (dB) & Low-EOG retention \\
\midrule
Zero anchor & 0.277 & \textbf{0.921} \\
Population calibration & 0.909 & 0.773 \\
Matched calibration & \textbf{2.463} & 0.843 \\
\bottomrule
\end{tabular}
\end{table}

The calibration controls separated reliable estimation from participant matching. Replacing the participant's propagation matrix with another participant's matrix degraded the reconstruction, while temporal shuffling reduced the usefulness of the EOG-derived guide. The stability criterion selected population calibration in four development evaluation cells. It did not, however, identify the owner of a stable matrix: the reliability-gated mismatched condition remained 0.0384 RRMSE worse than population calibration (0.0089--0.0697). The method therefore assumes that the calibration and evaluation segments belong to the same participant--session--task cell, as in the intended acquisition protocol. Figure~\ref{fig:natural-results} is designed to place the natural attenuation--retention trade-off beside these calibration controls.

\begin{figure*}[t]
  \centering
  \figureplaceholder{54mm}{\textbf{Natural-signal behavior and calibration controls.} Panel A should plot the aggregate operating points with low-EOG retention on the horizontal axis and EOG attenuation on the vertical axis for zero anchor, population calibration, matched calibration, mismatched calibration, shuffled EOG, and an evaluator-only upper bound. Panel B should show participant-level matched-minus-population effects for attenuation, retention, and EOG--EEG coherence. Panel C may show a scalp topography and power spectral density only if the corresponding evaluator arrays are available; use the lowest-identifier participant with complete outputs, the earliest eligible event, fixed channels, and identical scales across variants, and list those identifiers in the final caption. Panel D should show matched, population, and mismatched calibration controls without presenting the evaluator-only bound as an evaluation-time method.}
  \caption{Planned visualization of natural-recording attenuation, preservation, and calibration specificity. Natural analyses quantify relative behavior without assuming a clean natural target.}
  \Description{Placeholder for natural attenuation-retention plots, participant effects, and real spectral and topographic examples.}
  \label{fig:natural-results}
\end{figure*}

\subsection{Calibration representations and predictive intervals}

The auxiliary conditioning experiment tested whether the calibration benefit depended on the information conveyed rather than on the explicit matrix interface. The complete 15-participant analysis was inconclusive: its mean effect was $-1.4276$ ($-4.4154$--$0.0760$). One participant's entire support episode bank had an injected-artifact-to-clean-signal ratio of 167.2, compared with 0.33--4.02 for the other 14 participants. In the subsequent $[0.1,20]$ plausibility sensitivity analysis, a participant-matched calibration representation reduced RRMSE by 0.06083 relative to the zero-embedding reference (0.04064--0.08349; 14/14 participants). A mismatched representation increased error, giving matched-minus-mismatched separation of 0.16234 (0.13135--0.19460; 14/14). Matched-calibration gains remained positive for training pools from 30 to 259 participants, while the change from 30 to 259 was $-0.02236$ ($-0.04914$--$0.01582$), providing no evidence of an increasing gain over that range. Because the representation was fitted with paired clean calibration targets, this study is an evaluator-only upper bound that is consistent with, but does not independently validate, the calibration mechanism used by the main method.

The raw $K=32$ diffusion samples were too narrow: empirical 50\%, 80\%, and 90\% coverages were 0.302, 0.479, and 0.558, with CRPS 0.15162. Fold-held-out temperature scaling increased these values to 0.625, 0.802, and 0.851, with CRPS 0.15238. Combining propagation-matrix uncertainty with temperature scaling produced coverages of 0.622, 0.802, and 0.853 and CRPS 0.15031. The 80\% and 90\% intervals were therefore within 0.05 of their nominal levels, while the CRPS was lower than both the raw and temperature-only variants. Risk--coverage area also decreased from 0.05200 for the raw sample distribution to 0.05043. The 50\% interval remained conservative, indicating that one scalar adjustment did not align every nominal level equally well. Table~\ref{tab:secondary-results} reports both analyses, and Figure~\ref{fig:calibration-uq} specifies their final visual presentation.

\begin{table*}[t]
\caption{Controlled calibration-representation and development predictive-interval analyses. The injected representation study used a compact deterministic network with 15 evaluation participants and a plausible-range analysis of 14; the interval study used the 15-participant development diffusion model with $K=32$.}
\label{tab:secondary-results}
\centering
\small
\begin{tabularx}{\textwidth}{@{}p{.24\textwidth}Yp{.25\textwidth}p{.20\textwidth}@{}}
\toprule
Analysis & Comparison & Effect or interval & Support / interpretation \\
\midrule
Calibration representation & Matched vs zero embedding & 0.06083 (0.04064--0.08349) & 14/14 positive \\
Calibration specificity & Matched vs mismatched & 0.16234 (0.13135--0.19460) & 14/14 positive \\
Training-pool comparison & Gain at 259 minus gain at 30 & $-0.02236$ ($-0.04914$--$0.01582$) & No increasing trend detected \\
\midrule
Interval method & Nominal levels & Coverage at 50/80/90\% & CRPS \\
Raw sample distribution & 50/80/90\% & 0.302/0.479/0.558 & 0.15162 \\
Temperature scaling & 50/80/90\% & 0.625/0.802/0.851 & 0.15238 \\
Propagation uncertainty + temperature & 50/80/90\% & 0.622/0.802/0.853 & \textbf{0.15031} \\
\bottomrule
\end{tabularx}
\end{table*}

\begin{figure*}[t]
  \centering
  \figureplaceholder{54mm}{\textbf{Calibration representation and uncertainty.} Panel A should plot matched-versus-zero-embedding gain against training-pool size (30, 60, 120, 200, and 259 participants) with participant-bootstrap intervals. Panel B should show the matched, zero-embedding, and mismatched representation contrasts at the largest training pool. Label these panels as a controlled injected deterministic study with 15 evaluation participants and a plausible-range analysis of 14. Panel C should plot nominal versus empirical interval coverage for the raw sample distribution, temperature scaling, and propagation-uncertainty plus temperature scaling, with the identity line. Panel D should compare CRPS and risk--coverage area for the same interval variants. Use only stored aggregate values; do not invent sample bands or embedding projections.}
  \caption{Planned calibration-representation and predictive-interval figure. The left panels summarize a controlled injected deterministic study; the right panels evaluate sample-based predictive intervals on the 15-participant development data. Table~\ref{tab:secondary-results} reports the corresponding values.}
  \Description{Placeholder for calibration gain across training sizes, matched and mismatched conditions, interval coverage, CRPS, and risk-coverage area.}
  \label{fig:calibration-uq}
\end{figure*}

\section{Conclusion}
\label{sec:conclusion}

We introduced a subject-calibrated EOG-guided diffusion method for ocular-artifact removal. The method estimates an EOG-to-EEG propagation matrix from a separate calibration segment, shrinks uncertain individual estimates toward population calibration, and uses the resulting ocular contribution to condition a shared diffusion denoiser. Matched calibration improved paired waveform reconstruction in every development participant and in seven of eight held-out participants, with similar mean effect sizes in the two cohorts. Natural recordings showed a complementary gain in EOG attenuation and low-EOG signal retention. A controlled injected sensitivity analysis was consistent with participant matching when calibration information was delivered through a support-fitted representation, and propagation-aware sample widths improved the empirical coverage of predictive intervals on development data.

The present scope is EOG-assisted, same-session ocular-artifact removal with a compatible electrode montage. Natural recordings provide attenuation and preservation measures rather than a clean waveform target, and the predictive-interval analysis remains a development result. Linear regression and deterministic networks were competitive point estimators, so the current evidence supports the complete calibration-guided system rather than a general point-error advantage of diffusion. The next decisive evaluations are an external EEG--EOG cohort, a compute-matched comparison between diffusion sampling and deterministic ensembles, and held-out assessment of the predictive intervals. Within the evaluated setting, the results show that participant adaptation can be achieved by calibrating the ocular propagation pathway while keeping the EEG model shared.


\balance
\bibliographystyle{ACM-Reference-Format}
\bibliography{related_refs,references}

\end{document}
